\chapter*{Kết luận và Hướng phát triển}
\addcontentsline{toc}{chapter}{Kết luận và Hướng phát triển}

\section*{Kết luận}
\addcontentsline{toc}{section}{Kết luận}

Đồ án này đã giải quyết thành công bài toán tối ưu hóa chi phí Gas trên mạng lưới Ethereum cho các hệ thống Relayer/L2 Batching thông qua việc ứng dụng Offline Reinforcement Learning. Đi từ những thách thức cốt lõi của một môi trường biến động không ngừng (non-stationary) với yêu cầu khắt khe về thời gian hoàn thành (Zero-Miss SLA), chúng tôi đã rút ra được những kết luận mang tính định lượng và định tính quan trọng:

\begin{enumerate}
    \item \textbf{Thất bại của thuật toán tiêu chuẩn và Reward Shaping:} Các thuật toán tiên tiến như Decision Transformer (DT) \cite{chen2021decision} thất bại do mắc phải \textit{Majority Bias} khi học trên tập dữ liệu chuyên gia có độ lệch chuẩn cao (70\% thời gian giữ trạng thái chờ). Các thuật toán Q-Learning tiêu chuẩn cũng sụp đổ do \textit{Bất đối xứng Rủi ro - Lợi nhuận} (tỷ lệ 1:223,000) khiến Agent bị phạt quá nặng và chọn cách tích trữ an toàn. Kỹ thuật Reward Shaping \cite{ng1999policy} (giảm Penalty xuống mức hợp lý) là chìa khóa để Agent bắt đầu học được chiến lược "mua ở đáy".
    
    \item \textbf{Bất đối xứng Rủi ro - Lợi nhuận (Risk-Reward Asymmetry):} Thông qua phân tích toán học, chúng tôi định lượng tỷ lệ giữa hình phạt trễ hạn và lợi nhuận tiết kiệm là \textbf{223,000 : 1}. Đây là rào cản lớn nhất khiến các thuật toán Offline RL tiêu chuẩn bị rơi vào trạng thái "bi quan cực đoan" (Pessimistic Collapse). Việc tinh chỉnh Reward Shaping để cân bằng lại các thành phần này là điều kiện tiên quyết để mô hình có thể học được.
    
    \item \textbf{Kiến trúc Hệ thống Lai (Safe RL) là chuẩn mực thực tiễn:} Mặc dù mô hình AI (Discrete CQL) có khả năng tối ưu rất tốt, việc bàn giao 100\% quyền quyết định cho Neural Network trong môi trường tài chính rủi ro là không khả thi. Giải pháp đưa ra là Hệ thống Lai (Hybrid System) - giao cho AI 80\% thời gian đầu để "lướt sóng" tối ưu chi phí, và dùng một \textbf{Safety Layer} cứng (Rule-based) để ép bán ở 20\% thời gian cuối. Điều này giúp hạ Miss Rate từ 0.4\% xuống đúng 0.0\%.
    
    \item \textbf{Hiệu năng Vượt trội:} Trên tập dữ liệu đánh giá độc lập (Hold-out test) gồm 242 episodes đại diện cho 20\% thời gian cuối, mô hình lai của chúng tôi đã tiết kiệm được \textbf{11.1\% tổng chi phí Gas} so với chiến lược Baseline (Greedy) mà vẫn duy trì \textbf{0.0\% Tỷ lệ Trễ hạn (Miss Rate)}. Con số này chứng tỏ hệ thống đã khai thác được 28.6\% tiềm năng tuyệt đối của bộ giải Hindsight Oracle trong điều kiện không có thông tin tương lai.
\end{enumerate}

\section*{Hướng phát triển tương lai}
\addcontentsline{toc}{section}{Hướng phát triển tương lai}

Mặc dù đã đạt được những kết quả rất tích cực, nghiên cứu này vẫn còn những giới hạn mở ra cơ hội cho các nghiên cứu tiếp theo:

\begin{itemize}
    \item \textbf{Continuous Action Space và MILP Oracle:} Chuyển đổi từ không gian hành động rời rạc sang liên tục $[0, 1]$ kết hợp với các bộ giải Quy hoạch nguyên hỗn hợp (MILP) để tối ưu hóa chính xác đến từng đơn vị Wei, từ đó thu hẹp khoảng cách với trần lý thuyết.
    \item \textbf{Meta-Learning cho thị trường Non-stationary:} Ứng dụng các kiến trúc học thích nghi để mô hình tự động nhận diện và điều chỉnh chính sách theo từng trạng thái thị trường (Regime shift), thay vì sử dụng một bộ siêu tham số cố định.
    \item \textbf{Constrained Reinforcement Learning:} Tích hợp trực tiếp các ràng buộc an toàn (Safety Constraints) vào hàm mục tiêu thông qua các phương pháp Lagrangian hoặc CPO (Constrained Policy Optimization), từ đó loại bỏ sự phụ thuộc vào lớp Safety Layer cứng bên ngoài.
\end{itemize}
